% Imporditud: tools/import_eksperimendid.py
% Allikas: Eesti füüsikaolümpiaadi eksperimendi ülesannete
%   kogu (Rael Kalda, 2023), ülesanne Ü181, lahendus L181.
\setAuthor{Jaan Kalda}
\setRound{lõppvoor}
\setYear{2022}
\setNumber{G E2}
\setDifficulty{5}
\setTopic{dünaamika}
\setVahendid{kaks pingpongi palli, millest ühe külge on läbi nõelaaugu kinnitatud niit; puitjoonlaud; plastjoonlaud}
\setPunktid{14}
\setAllikas{Eksperimendi kogu Ü181, lahendus lk 143}
\setLahendusseis{toores}

\prob{Pingpong}
Leidke pingpongi palli ja puitjoonlaua vaheline hõõrdetegur.

\hint

\solu
Ehitame puitjoonlauast kaldpinna toetades ühe otsa vastu lauapinda ja kergitades

teist käega. Hoiame samal ajal samas käes vertikaalselt plastjoonlauda niimoodi,

et saame selle abil mõõta puitjoonlaua otsa kõrgust h lauapinnast. Seda kõrgust

kasutame selleks, et teha kindlaks kaldpinna nurga $\alpha$ = arcsin h/L, kus L on puit-

joonalua kogupikkus. Leiame sellise joonlaua kaldenurga, mille puhul on võimalik

hoida pingpongi palli kaldpinna peal paigal tõmmates niidist, mis on paralleelne

kaldpinnaga, kusjuures niidi kinnituspunkt on palli puutepunktist diameetri kau-

gusel, vt joonist. Paralleelsuse tagamiseks on mugav kasutada teist pingpongi palli,

mis hoiab niiti samal kaugusel joonlauast. Tasakaalutingimus on, et kolm raken-

datud jõudu --- raskusjõud, hõõrdejõu ja toereaktsiooni resultant (mille siht moo-

dustab nurga $\beta$ = arctan $\mu$ pinnanormaaliga) ja niidi pinge lõikuvad ühes punktis

P. On lihtne näha, et tan $\beta$ = 12 tan $\alpha$, st $\mu$ = 1 tan $\alpha$.

2b

TP a N+Fh b

mg

a
\probend
